\chapter{Celiac Disease Serology}

\section*{What Is This Test?}
Celiac disease serology detects antibodies produced in response to gluten ingestion in genetically susceptible individuals. Celiac disease is an immune-mediated enteropathy triggered by gluten (the storage protein in wheat, barley, and rye) in people carrying the HLA-DQ2 or HLA-DQ8 haplotypes \cite{lebwohl2018coeliac}. The antibody response is directed against tissue transglutaminase (tTG), the enzyme that deamidates gluten peptides, and against deamidated gliadin peptides (DGP) and endomysium. The first-line test is the anti-tissue transglutaminase IgA antibody (tTG-IgA), which has high sensitivity and specificity in untreated celiac disease, along with a total IgA level to exclude IgA deficiency (which causes false-negative IgA-based tests) \cite{rubiotapia2013acg}. In IgA-deficient patients, tTG-IgG and DGP-IgG are used. Anti-endomysial antibody (EMA-IgA), a more specific but less widely automated test, is used to confirm borderline positive results. Positive serology must be confirmed by duodenal biopsy for a definitive diagnosis of celiac disease in most patients, though 2012 European Society for Paediatric Gastroenterology, Hepatology and Nutrition (ESPGHAN) guidelines permit a no-biopsy diagnosis in symptomatic children with tTG-IgA greater than 10 times the upper limit of normal and a positive EMA \cite{husby2012espghan}. The test is only valid in patients eating gluten, because antibodies decline when gluten is removed.

\section*{Alternative Names}
\begin{itemize}
  \item Celiac Panel
  \item Celiac Blood Test
  \item Tissue Transglutaminase Antibody (tTG-IgA, tTG-IgG)
  \item Anti-Endomysial Antibody (EMA)
  \item Deamidated Gliadin Peptide Antibody (DGP-IgA, DGP-IgG)
  \item Celiac Disease Antibody Panel
\end{itemize}
\section*{Principle}
Celiac serology uses enzyme-linked immunosorbent assay (ELISA) or automated chemiluminescent immunoassay (CLIA) to detect IgA and IgG antibodies against tissue transglutaminase and deamidated gliadin peptides. The antigen (recombinant human tTG, or deamidated gliadin peptides) is immobilized on the solid phase; patient serum is added, and bound antibody is detected with an enzyme- or label-conjugated anti-human antibody, producing a signal proportional to antibody concentration. Total IgA is measured simultaneously (by nephelometry or immunoassay) because IgA deficiency, present in 1:300--1:500 people and more common in celiac disease, causes false-negative tTG-IgA and EMA-IgA results. Anti-endomysial antibody is detected by indirect immunofluorescence on monkey esophagus or human umbilical cord sections, where a characteristic reticulin-like staining pattern of the endomysium is seen. In IgA-deficient patients (total IgA below the age-specific lower limit), tTG-IgG and DGP-IgG are the appropriate tests. HLA-DQ2/DQ8 typing, when used, is performed by PCR-based methods on DNA and has a very high negative predictive value (almost no one without DQ2/DQ8 develops celiac disease).

\section*{History}
The link between gluten and celiac disease was established during the Dutch famine of World War II, when pediatrician Willem Dicke observed that children with celiac disease improved dramatically when wheat was unavailable and relapsed when it was reintroduced; his 1950 thesis described this. The small intestinal lesion (villous atrophy) was characterized by John Paul in 1954. The autoantibody basis emerged in the 1960s--1980s: reticulin antibodies (1971), gliadin antibodies (1980s), and anti-endomysial antibodies (1983). In 1997, Dieterich and colleagues identified tissue transglutaminase as the autoantigen of endomysial antibodies \cite{dieterich1997identification}, enabling the development of the quantitative tTG-IgA assay that remains the first-line test. Deamidated gliadin peptide (DGP) assays were introduced in the 2000s and improved the serologic diagnosis of celiac disease, particularly in IgA deficiency. The discovery of the HLA-DQ2/DQ8 association in the 1980s and 1990s clarified genetic susceptibility.

\section*{Why This Test Is Ordered}
\begin{itemize}
  \item Evaluation of chronic or recurrent diarrhea, bloating, abdominal pain, and weight loss
  \item Investigation of iron deficiency anemia that is unexplained or refractory to oral iron
  \item Evaluation of malabsorption with low ferritin, folate, vitamin B12, or fat-soluble vitamins
  \item Assessment of patients with unexplained fatigue, osteoporosis or early-onset osteopenia, or dental enamel defects
  \item Screening of first-degree relatives of patients with celiac disease (approximately 10\% risk) and of patients with type 1 diabetes, autoimmune thyroid disease, Down syndrome, Turner syndrome, or IgA deficiency
  \item Evaluation of unexplained transaminitis (elevated ALT/AST)
  \item Investigation of dermatitis herpetiformis (an intensely itchy vesicular rash, the cutaneous manifestation of celiac disease), which may be present without gastrointestinal symptoms
  \item Workup of infertility, recurrent miscarriage, or unexplained short stature in children
  \item Monitoring adherence to a gluten-free diet (falling antibody titers)
  \item Evaluation of persistent symptoms in patients already on a gluten-free diet who were never formally diagnosed
\end{itemize}

\section*{Primary vs Secondary Test}
tTG-IgA with total IgA is a primary screening test for celiac disease. It is not recommended in asymptomatic, average-risk people without risk factors. Positive or borderline serology requires confirmatory testing: EMA-IgA for borderline cases and, in most guidelines, duodenal biopsy for a definitive diagnosis. In children with very high tTG-IgA (greater than 10 times the upper limit of normal) plus positive EMA, biopsy may be omitted per ESPGHAN 2012 guidelines.

\section*{How This Test Is Performed}
A blood sample is collected by standard venipuncture into a serum separator tube. The panel typically measures tTG-IgA, total IgA, and (reflexively, if IgA is deficient) tTG-IgG and DGP-IgG. EMA-IgA is performed by indirect immunofluorescence when needed. The patient MUST be eating a gluten-containing diet at the time of testing for the result to be valid. Results are usually available within 1--3 days for the automated tTG assay and 3--7 days for EMA if it must be sent out. If a patient has already been on a gluten-free diet, the test may be falsely negative; serology is not useful in this setting, and a gluten challenge (reintroduction of gluten) or a biopsy with specific HLA testing may be required.

\section*{What Preparation Is Needed}
No fasting is required, but the patient must be on a normal gluten-containing diet (at least one serving of gluten daily for 6--12 weeks before testing) for the serology to be positive in untreated disease. Patients who have already started a gluten-free diet should not stop and restart gluten abruptly without medical supervision, because a gluten challenge can cause severe symptoms; the clinical team should guide the timing and duration. Medications such as immunosuppressants (including corticosteroids) can lower antibody levels and cause false-negative results.

\section*{Contraindications and Precautions}
\begin{itemize}
  \item There are no absolute contraindications to standard venipuncture. Important interpretive cautions:
  \item (1) IgA deficiency (present in 1:300--1:500 of the general population and 2--3\% of celiac patients) causes false-negative tTG-IgA and EMA-IgA; total IgA must always be measured, and IgG-based tests used if deficient.
  \item (2) The test is only valid on a gluten-containing diet; a gluten-free diet markedly lowers or normalizes antibody levels.
  \item (3) False-positive tTG-IgA is reported in autoimmune hepatitis, cirrhosis, primary biliary cholangitis, inflammatory bowel disease, heart failure, and some infections; EMA is more specific than tTG and can help clarify borderline results.
  \item (4) Immunosuppression (including steroids and immunosuppressant biologics) can lower antibody titers and produce false negatives.
  \item (5) In children younger than 2 years, tTG and EMA are less reliable, and DGP-IgA/IgG is more sensitive.
  \item (6) Antibody titers decline with treatment but may take 6--12 months to normalize; rising titers in a treated patient suggest dietary non-adherence or (rarely) complications such as refractory celiac disease.
  \item (7) Serology does not establish the diagnosis alone in adults; villous atrophy on duodenal biopsy (Marsh 3) is the diagnostic gold standard in most guidelines \cite{altoma2019european}.
\end{itemize}
\section*{Risks}
Minimal risk. Slight pain or bruising at the venipuncture site. Rarely: hematoma, infection, or vasovagal syncope.

\section*{What Happens After the Test}
The result is interpreted in the context of symptoms and total IgA. A strongly positive tTG-IgA (greater than 10 times the upper limit of normal) with a positive EMA supports the diagnosis; in symptomatic children, this may be sufficient for diagnosis without biopsy. In most adults and in borderline cases, upper gastrointestinal endoscopy with duodenal biopsies is performed to confirm villous atrophy. A negative result in a patient on a gluten-containing diet makes celiac disease very unlikely, though the condition is not entirely excluded in IgA deficiency (if IgG testing was not performed) or in patients already on a gluten-free diet. Patients diagnosed with celiac disease are referred to a dietitian and started on a strict gluten-free diet, with serologic monitoring (tTG-IgA) at 3--6 months and 12 months and then annually to assess adherence. Patients should also be screened for associated conditions (thyroid disease, type 1 diabetes, osteoporosis) and counseled on the importance of lifelong gluten avoidance.

\section*{Understanding Results}

\subsection*{Below Normal (Negative)}
\begin{itemize}
  \item \textbf{Negative tTG-IgA on a gluten-containing diet:} Celiac disease is very unlikely (negative predictive value greater than 95\% for symptomatic patients). If clinical suspicion remains high (e.g., refractory iron deficiency anemia or malabsorption), duodenal biopsy or HLA-DQ2/DQ8 typing can be considered.
  \item \textbf{Negative result with IgA deficiency:} The tTG-IgA result is uninterpretable; tTG-IgG and DGP-IgG must be checked, because celiac disease is more common in IgA deficiency.
  \item \textbf{Negative result in a patient already on a gluten-free diet:} The result is invalid for diagnosis; antibody levels normalize on gluten exclusion.
\end{itemize}

\subsection*{Above Normal (Positive)}
\begin{itemize}
  \item \textbf{Positive tTG-IgA:} Highly suggestive of celiac disease when the patient is on a gluten-containing diet; titers greater than 10 times the upper limit of normal are strongly associated with Marsh 3 villous atrophy. Confirmatory EMA and/or duodenal biopsy are performed per guidelines.
  \item \textbf{Borderline tTG-IgA (1--3 times the upper limit of normal):} Repeat testing, EMA confirmation, and biopsy are warranted because false positives occur (autoimmune liver disease, inflammatory bowel disease, infections, congestive heart failure).
  \item \textbf{Positive EMA-IgA:} High specificity for celiac disease; used to confirm a positive or borderline tTG result.
  \item \textbf{Positive tTG-IgG or DGP-IgG with IgA deficiency:} Consistent with celiac disease in IgA-deficient patients.
  \item \textbf{Serologic resolution of titers on a gluten-free diet:} Confirms the diagnosis in patients who had villous atrophy and indicates good dietary adherence when titers normalize.
\end{itemize}

\section*{Normal Results}
\begin{itemize}
  \item \textbf{tTG-IgA:} Negative or below the assay cutoff (typically less than 20 U/mL, laboratory-dependent); positive results are reported with the degree of elevation
  \item \textbf{EMA-IgA:} Negative
  \item \textbf{DGP-IgA and DGP-IgG:} Negative (used mainly in IgA deficiency and in children under 2 years)
  \item \textbf{Total IgA:} Within the age-specific reference range; low total IgA indicates deficiency
  \item \textbf{Reference note:} Reference cutoffs vary between manufacturers and laboratories; serial testing should use the same laboratory because titers are not interchangeable across assays
\end{itemize}

\section*{Follow-up Tests}
\begin{itemize}
  \item \textbf{Total IgA:} Always measured with the tTG-IgA to interpret the IgA-based result
  \item \textbf{Anti-endomysial antibody (EMA-IgA):} Confirmatory, highly specific test for borderline or unexpected positive results
  \item \textbf{tTG-IgG and DGP-IgG:} First-line tests in IgA-deficient patients
  \item \textbf{Upper endoscopy with duodenal biopsy:} The diagnostic gold standard in most adults and in borderline pediatric cases; multiple biopsies from the duodenal bulb and distal duodenum are evaluated for villous atrophy (Marsh grading)
  \item \textbf{HLA-DQ2/DQ8 typing:} Not diagnostic alone, but very high negative predictive value; useful to exclude celiac disease in patients already on a gluten-free diet and in first-degree relatives with ambiguous serology
  \item \textbf{Iron studies, folate, vitamin B12, fat-soluble vitamins:} Screen for coexistent deficiency from malabsorption
  \item \textbf{Bone density testing (DXA):} Because celiac disease is associated with osteoporosis and osteomalacia from calcium and vitamin D malabsorption
  \item \textbf{Thyroid function tests and thyroid antibodies:} Autoimmune thyroid disease is associated with celiac disease
  \item \textbf{Type 1 diabetes screening:} In appropriate clinical context, because of the strong association
  \item \textbf{Serial tTG-IgA on a gluten-free diet:} Monitor adherence; titers fall within 3--6 months and should normalize within 12 months
\end{itemize}

\section*{References}
\bibliographystyle{unsrtnat}
\bibliography{references}

\clearpage
\section*{Regulatory Status and Authorization Date}
This chapter describes a test category, not a specific commercial product. FDA, CDSCO, EU IVDR, and MHRA approval or registration dates therefore cannot be assigned without the assay manufacturer, product name, instrument, and intended use. For laboratory-developed tests, an individual agency approval date may not exist. See the Preface for the interpretation of regulatory status.

\clearpage
